\section{常微分方程周期性}

% Source: Google Drive file 1D4SwMoS2wPcTdw3u-nUG55nyX6Pa2w3y
% Grade-2 / 大二 notes. Energy + uniqueness argument for a conservative oscillator,
% originally written for a Yau 2020-style problem \(\ddot x+x+x^3=0\).
% The same argument is the closed-orbit half of CAM 2026 Spring Problem 6.

There exists a constant \(C\) such that
\[
(x')^{2}+x^{2}+\tfrac12 x^{4}=C.
\]
Then \(x\) is bounded, since
\[
0\le x^{2}+\tfrac12 x^{4}\le C,
\]
hence \(\lvert x\rvert<M\) for some \(M>0\).

If for every \(t_{0}>0\) one has \(x'(t)\neq 0\) for all \(t>t_{0}\), then \(x\) is eventually monotonic, so \(\lim_{t\to+\infty}x(t)=s\in\mathbb{R}\). There exist sequences \(t_{n}\to\infty\) and \(h_{n}\to\infty\) with
\[
\lim_{n\to\infty}x'(t_{n})=\lim_{n\to\infty}x''(h_{n})=0.
\]
The energy identity and the ODE then give \(s^{2}+\tfrac12 s^{4}=C\) and \(s+s^{3}=0\), hence \(s=0\), \(C=0\), and \(x\equiv 0\).

Otherwise there is a sequence \(t_{n}\to\infty\) with \(x'(t_{n})=0\). Then \(x^{2}(t_{n})+\tfrac12 x^{4}(t_{n})=C\) has only finitely many roots, so (passing to a subsequence) one finds \(t_{1}\neq t_{2}\) with
\[
x(t_{1})=x(t_{2}),\qquad x'(t_{1})=x'(t_{2})=0.
\]
The shifted solution \(x(t-t_{1}+t_{2})\) and \(x(t)\) share the same Cauchy data at \(t_{1}\). Uniqueness for the Lipschitz ODE \(x''=-x-x^{3}\) forces them to coincide, so \(x\) is periodic of period \(t_{2}-t_{1}\).
